%! Potential Errors When Performing Tests — Unit 6, Lesson 6.7 — Homework (Answer Key)
\documentclass[10pt]{article}
\usepackage{apstats-article}
\usepackage{apstats-key}   % pulls in -boxes; do NOT also load -boxes
\usepackage{grading}
\gradingassignment{apstats-u06-l07-hw}
\gradingcourse{advanced_statistics}
\newcommand{\TallMath}[1]{$\displaystyle #1\rule[-1.4em]{0pt}{3.2em}$}

\begin{document}

\pageheader{Unit 6, Lesson 6.7}{Homework --- Answer Key}
{\small\itshape Write each final answer in its box. Only what is inside a box is graded.}\par

\vspace{0.15in}

\begin{enumerate}

  \item Blood pressure test:
        \begin{enumerate}[label=\alph*.]
          \item \ansline{Type I error: concluding that more than 10\% of adults in the city have high blood pressure when the true rate is actually 10\%. The health department would allocate resources for a problem that doesn't exist.}\gradedkeyline{\gradedblank[id=hw1a, type=rubric, points=2, display={Concluding more than 10\% of the city's adults have high blood pressure when the true rate is 10\%.}]{\linewidth+6pt}{\gradedlinesheight{1}}}
          \item \ansline{Type II error: failing to detect that more than 10\% of adults have high blood pressure when the true rate exceeds 10\%. The department would miss a public health problem.}\gradedkeyline{\gradedblank[id=hw1b, type=rubric, points=2, display={Failing to conclude more than 10\% have high blood pressure when the true rate really exceeds 10\%.}]{\linewidth+6pt}{\gradedlinesheight{1}}}
          \item \ansline{Power would be higher when $p=0.20$ than when $p=0.12$ because the true proportion is farther from $p_0=0.10$, making the effect easier to detect.}\gradedkeyline{\gradedblank[id=hw1c, type=rubric, points=2, display={Lower at 12\%: 0.12 is closer to $p_0 = 0.10$ than 0.20 is, so the difference is harder to detect.}]{\linewidth+6pt}{\gradedlinesheight{1}}}
        \end{enumerate}

  \newpage
  \item Reading intervention:
        \begin{enumerate}[label=\alph*.]
          \item \ansline{From the principal's perspective, a Type II error (missing a real benefit) is more costly --- a genuinely effective intervention would not be implemented, and students would be denied an improvement in their education. A Type I error means implementing an ineffective program, wasting resources but not directly harming students.}\gradedkeyline{\gradedblank[id=hw2a, type=rubric, points=2, display={Type II: an intervention that really works would not be adopted, so students lose a real benefit (a Type I error only wastes resources).}]{\linewidth+6pt}{\gradedlinesheight{2}}}
          \item \ansline{Increasing from $n=60$ to $n=240$ increases power (reduces $SE$) and decreases $\beta$ (probability of Type II error), making the test more sensitive to a real improvement.}\gradedkeyline{\gradedblank[id=hw2b, type=rubric, points=2, display={Power increases and $P(\text{Type II error})$ decreases (larger $n$, smaller $SE$).}]{\linewidth+6pt}{\gradedlinesheight{1}}}
        \end{enumerate}

  \newpage
  \item True/False:
        \begin{enumerate}[label=\alph*.]
          \item \ans{F} \kern3pt\gradedtop{\gradedblank[id=hw3a1, type=choice, points=1, answer=F]{1cm}{0.17in}}\kern0.6em\ansline{Reducing $\alpha$ decreases power (makes it harder to reject $H_0$), increasing $\beta$.}\gradedkeyline{\gradedblank[id=hw3a2, type=rubric, points=1, display={False: a smaller $\alpha$ makes $H_0$ harder to reject, so power decreases.}]{\linewidth+6pt}{\gradedlinesheight{1}}}
          \item \ans{F} \kern3pt\gradedtop{\gradedblank[id=hw3b1, type=choice, points=1, answer=F]{1cm}{0.17in}}\kern0.6em\ansline{A Type II error occurs when we \emph{fail to} reject $H_0$. Specifically, it happens when $H_0$ is NOT rejected but $H_a$ is actually true.}\gradedkeyline{\gradedblank[id=hw3b2, type=rubric, points=1, display={False: a Type II error happens when we FAIL to reject $H_0$ while $H_a$ is true.}]{\linewidth+6pt}{\gradedlinesheight{1}}}
        \end{enumerate}

\end{enumerate}

\begin{reflectionbox}
\textbf{Preview:} Lesson 6.8 introduces the two-sample $z$-interval for $p_1 - p_2$.
\end{reflectionbox}

\end{document}
